%%%%%%%% ICML 2025 EXAMPLE LATEX SUBMISSION FILE %%%%%%%%%%%%%%%%%

\documentclass{article}

% Recommended, but optional, packages for figures and better typesetting:
\usepackage{microtype}
\usepackage{graphicx}
\usepackage{subfigure}
\usepackage{booktabs} % for professional tables

% hyperref makes hyperlinks in the resulting PDF.
% If your build breaks (sometimes temporarily if a hyperlink spans a page)
% please comment out the following usepackage line and replace
% \usepackage{icml2025} with \usepackage[nohyperref]{icml2025} above.
\usepackage{hyperref}


% Attempt to make hyperref and algorithmic work together better:
\newcommand{\theHalgorithm}{\arabic{algorithm}}

% Use the following line for the initial blind version submitted for review:
\usepackage{icml2025}

% If accepted, instead use the following line for the camera-ready submission:
% \usepackage[accepted]{icml2025}

% For theorems and such
\usepackage{amsmath}
\usepackage{amssymb}
\usepackage{mathtools}
\usepackage{amsthm}

% if you use cleveref..
\usepackage[capitalize,noabbrev]{cleveref}

%%%%%%%%%%%%%%%%%%%%%%%%%%%%%%%%
% THEOREMS
%%%%%%%%%%%%%%%%%%%%%%%%%%%%%%%%
\theoremstyle{plain}
\newtheorem{theorem}{Theorem}[section]
\newtheorem{proposition}[theorem]{Proposition}
\newtheorem{lemma}[theorem]{Lemma}
\newtheorem{corollary}[theorem]{Corollary}
\theoremstyle{definition}
\newtheorem{definition}[theorem]{Definition}
\newtheorem{assumption}[theorem]{Assumption}
\theoremstyle{remark}
\newtheorem{remark}[theorem]{Remark}

% Todonotes is useful during development; simply uncomment the next line
%    and comment out the line below the next line to turn off comments
%\usepackage[disable,textsize=tiny]{todonotes}
\usepackage[textsize=tiny]{todonotes}


% The \icmltitle you define below is probably too long as a header.
% Therefore, a short form for the running title is supplied here:
\icmltitlerunning{Submission and Formatting Instructions for ICML 2025}

\begin{document}


\onecolumn
% \twocolumn[
\icmltitle{Submission and Formatting Instructions for \\
           International Conference on Machine Learning (ICML 2025)}

% It is OKAY to include author information, even for blind
% submissions: the style file will automatically remove it for you
% unless you've provided the [accepted] option to the icml2025
% package.

% List of affiliations: The first argument should be a (short)
% identifier you will use later to specify author affiliations
% Academic affiliations should list Department, University, City, Region, Country
% Industry affiliations should list Company, City, Region, Country

% You can specify symbols, otherwise they are numbered in order.
% Ideally, you should not use this facility. Affiliations will be numbered
% in order of appearance and this is the preferred way.
\icmlsetsymbol{equal}{*}

\begin{icmlauthorlist}
\icmlauthor{Firstname1 Lastname1}{equal,yyy}
\icmlauthor{Firstname2 Lastname2}{equal,yyy,comp}
\icmlauthor{Firstname3 Lastname3}{comp}
% \icmlauthor{Firstname4 Lastname4}{sch}
% \icmlauthor{Firstname5 Lastname5}{yyy}
% \icmlauthor{Firstname6 Lastname6}{sch,yyy,comp}
% \icmlauthor{Firstname7 Lastname7}{comp}
% %\icmlauthor{}{sch}
% \icmlauthor{Firstname8 Lastname8}{sch}
% \icmlauthor{Firstname8 Lastname8}{yyy,comp}
%\icmlauthor{}{sch}
%\icmlauthor{}{sch}
\end{icmlauthorlist}

% \icmlaffiliation{yyy}{Department of XXX, University of YYY, Location, Country}
% \icmlaffiliation{comp}{Company Name, Location, Country}
% \icmlaffiliation{sch}{School of ZZZ, Institute of WWW, Location, Country}

% \icmlcorrespondingauthor{Firstname1 Lastname1}{first1.last1@xxx.edu}
% \icmlcorrespondingauthor{Firstname2 Lastname2}{first2.last2@www.uk}

% You may provide any keywords that you
% find helpful for describing your paper; these are used to populate
% the "keywords" metadata in the PDF but will not be shown in the document
\icmlkeywords{Machine Learning, ICML}

\vskip 0.3in
% ]


% this must go after the closing bracket ] following \twocolumn[ ...

% This command actually creates the footnote in the first column
% listing the affiliations and the copyright notice.
% The command takes one argument, which is text to display at the start of the footnote.
% The \icmlEqualContribution command is standard text for equal contribution.
% Remove it (just {}) if you do not need this facility.

%\printAffiliationsAndNotice{}  % leave blank if no need to mention equal contribution
% \printAffiliationsAndNoti ce{\icmlEqualContribution} % otherwise use the standard text.

\begin{abstract}
\end{abstract}

\input{tex/Introduction}

\section{Second-order Influence}
% =========================
% Second-order influence: derivation details
% =========================

\paragraph{Setup.}
Let the (average) training objective be
\[
\mathcal{L}(\theta) \;=\; \frac{1}{N}\sum_{z\in \mathcal{D}_{\mathrm{train}}}\ell(z,\theta),
\qquad
\mathcal{L}_S(\theta) \;=\; \frac{1}{N}\sum_{z\in S}\ell(z,\theta),
\]
and define the perturbed objective (upweighting the subset $S$ by $\varepsilon$)
\[
\mathcal{L}_\varepsilon(\theta)\;=\;\mathcal{L}(\theta)+\varepsilon\,\mathcal{L}_S(\theta).
\]
Let $\theta^\star=\arg\min_\theta \mathcal{L}(\theta)$ and $\theta_\varepsilon=\arg\min_\theta \mathcal{L}_\varepsilon(\theta)$.
Define
\[
g_f := \nabla_\theta f(\theta^\star),\quad H_f:=\nabla_\theta^2 f(\theta^\star),
\]
\[
g_S := \nabla_\theta \mathcal{L}_S(\theta^\star)=\frac{1}{N}\sum_{z\in S}\nabla_\theta \ell(z,\theta^\star),\quad
H_S := \nabla_\theta^2 \mathcal{L}_S(\theta^\star)=\frac{1}{N}\sum_{z\in S}\nabla_\theta^2 \ell(z,\theta^\star),
\]
\[
H := \nabla_\theta^2 \mathcal{L}(\theta^\star)=\frac{1}{N}\sum_{z\in \mathcal{D}_{\mathrm{train}}}\nabla_\theta^2 \ell(z,\theta^\star).
\]

\paragraph{Step 1: expansion of the stationarity condition.}
The optimality condition for $\theta_\varepsilon$ is
\[
\nabla_\theta \mathcal{L}_\varepsilon(\theta_\varepsilon)
=\nabla_\theta\mathcal{L}(\theta_\varepsilon)+\varepsilon\,\nabla_\theta\mathcal{L}_S(\theta_\varepsilon)=0.
\]
Write $\Delta := \theta_\varepsilon-\theta^\star$ (so $\Delta=O(\varepsilon)$), and use
\[
\nabla_\theta\mathcal{L}(\theta^\star+\Delta)\approx \nabla_\theta\mathcal{L}(\theta^\star)+H\Delta
=H\Delta
\quad(\text{since }\nabla_\theta\mathcal{L}(\theta^\star)=0),
\]
\[
\nabla_\theta\mathcal{L}_S(\theta^\star+\Delta)\approx g_S+H_S\Delta.
\]
Keeping terms up to order $O(\varepsilon^2)$ yields
\[
0 \approx H\Delta + \varepsilon(g_S+H_S\Delta)
\quad\Longrightarrow\quad
(H+\varepsilon H_S)\Delta \approx -\varepsilon g_S.
\]
Using the Neumann expansion
\[
(H+\varepsilon H_S)^{-1}\approx H^{-1}-\varepsilon H^{-1}H_S H^{-1},
\]
we obtain
\[
\Delta
\approx -(H^{-1}-\varepsilon H^{-1}H_S H^{-1})\,\varepsilon g_S
= -\varepsilon H^{-1}g_S + \varepsilon^2 H^{-1}H_S H^{-1}g_S
\;+\;O(\varepsilon^3).
\]

\paragraph{Step 2: second-order Taylor expansion of $f(\theta_\varepsilon)$.}
A second-order expansion around $\theta^\star$ gives
\[
f(\theta_\varepsilon) = f(\theta^\star+\Delta)
\approx f(\theta^\star) + g_f^\top \Delta + \frac{1}{2}\Delta^\top H_f \Delta.
\]
Substituting $\Delta = -\varepsilon H^{-1}g_S + \varepsilon^2 H^{-1}H_S H^{-1}g_S + O(\varepsilon^3)$
and keeping terms up to $O(\varepsilon^2)$:
\[
g_f^\top\Delta
= -\varepsilon\, g_f^\top H^{-1}g_S
+\varepsilon^2\, g_f^\top H^{-1}H_S H^{-1}g_S,
\]
\[
\frac{1}{2}\Delta^\top H_f \Delta
= \frac{\varepsilon^2}{2}\,(H^{-1}g_S)^\top H_f (H^{-1}g_S).
\]
Therefore,
\[
f(\theta_\varepsilon)
\approx f(\theta^\star)
-\varepsilon\, g_f^\top H^{-1}g_S
+\frac{\varepsilon^2}{2}\Big(
(H^{-1}g_S)^\top H_f (H^{-1}g_S)
+2(H^{-1}g_f)^\top H_S (H^{-1}g_S)
\Big),
\]
where we used $g_f^\top H^{-1}H_S H^{-1}g_S = (H^{-1}g_f)^\top H_S (H^{-1}g_S)$ (symmetry).

% =========================
% Pairwise expansion of the synergy term
% =========================

\paragraph{Synergy term as pairwise interactions.}
Define $v_i := H^{-1}\nabla_\theta \ell(z_i,\theta^\star)$. Then
\[
H^{-1}g_S
=\frac{1}{N}\sum_{z_i\in S} H^{-1}\nabla_\theta \ell(z_i,\theta^\star)
=\frac{1}{N}\sum_{z_i\in S} v_i.
\]
Hence,
\[
(H^{-1}g_S)^\top H_f(H^{-1}g_S)
=\Big(\frac{1}{N}\sum_{i\in S} v_i\Big)^\top
H_f
\Big(\frac{1}{N}\sum_{j\in S} v_j\Big)
=\frac{1}{N^2}\sum_{i\in S}\sum_{j\in S} v_i^\top H_f v_j.
\]
This shows that synergy is governed by pairwise similarities $v_i^\top H_f v_j$.

% =========================
% Cross-dataset decomposition
% =========================

\paragraph{Cross-dataset synergy (two groups).}
Let $S=S_D\cup S_H$ and $g_S=g_D+g_H$. Then
\[
(H^{-1}g_S)^\top H_f(H^{-1}g_S)
=
(H^{-1}g_D)^\top H_f(H^{-1}g_D)
+
(H^{-1}g_H)^\top H_f(H^{-1}g_H)
+
2(H^{-1}g_D)^\top H_f(H^{-1}g_H).
\]
Moreover, with $v_i:=H^{-1}\nabla_\theta \ell(z_i,\theta^\star)$,
\[
2(H^{-1}g_D)^\top H_f(H^{-1}g_H)
=\frac{2}{N^2}\sum_{i\in S_D}\sum_{j\in S_H} v_i^\top H_f v_j.
\]

% =========================
% Low-dimensional projection approximation
% =========================

\paragraph{Low-dimensional projection approximation.}
Let $P\in\mathbb{R}^{k\times n}$ be a projection (with $k\ll n$).
A common approximation for first-order influence is
\[
g_{\mathrm{te}}^\top H^{-1}g_{\mathrm{tr}}
\;\approx\;
(Pg_{\mathrm{te}})^\top (P H P^\top)^{-1}(Pg_{\mathrm{tr}}).
\]
For the synergy term, define the projected solve
\[
u_S := (P H P^\top)^{-1}(P g_S)\in\mathbb{R}^{k},
\]
so that $H^{-1}g_S \approx P^\top u_S$ in the projected subspace, and approximate
\[
(H^{-1}g_S)^\top H_f(H^{-1}g_S)
\;\approx\;
u_S^\top (P H_f P^\top) u_S
=
\big((P H P^\top)^{-1}P g_S\big)^\top (P H_f P^\top)\big((P H P^\top)^{-1}P g_S\big).
\]

% ============================================================
% Synthetic data selection that complements real data
% via second-order influence (LLM setting)
% ============================================================

\section{Selecting Synthetic Data that Complements Real Data via Second-Order Influence}
\label{sec:synth_complement}

\subsection{Motivation}
Large language models are increasingly tuned using a mixture of (i) small, high-quality real data
(human-written demonstrations or preferences) and (ii) large pools of synthetic data
(distilled or self-generated instructions, dialogues, preference pairs).
A key challenge is that \emph{the value of a synthetic example depends on what real data is already present}.
Synthetic data can be:
\begin{itemize}
  \item \textbf{redundant}: it pushes the parameters in directions already covered by real data, yielding diminishing returns;
  \item \textbf{complementary}: it fills gaps left by real data and improves target evaluation performance;
  \item \textbf{interfering}: it interacts with real data in a way that harms the target behavior (negative transfer).
\end{itemize}

First-order influence scores are linear in the chosen subset and thus cannot represent
interaction effects between real and synthetic groups.
Second-order influence introduces a quadratic term that captures these interactions,
providing a principled signal for \emph{complementary synthetic subset selection}.

\subsection{Training objective and data groups}
Let $z=(x,y)$ be a supervised instruction-following example.
Let $\mathcal{D}_{\mathrm{real}}$ denote the real dataset and $\mathcal{P}_{\mathrm{syn}}$ a large pool of synthetic candidates.
We aim to select a subset $S \subseteq \mathcal{P}_{\mathrm{syn}}$ under a budget constraint
(e.g., $|S|\le B$ or token budget $\mathrm{Tok}(S)\le T$).

Define the (average) training losses
\[
\mathcal{L}_{\mathrm{real}}(\theta) := \frac{1}{|\mathcal{D}_{\mathrm{real}}|}\sum_{z\in \mathcal{D}_{\mathrm{real}}} \ell(z,\theta),
\qquad
\mathcal{L}_{S}(\theta) := \frac{1}{|S|}\sum_{z\in S} \ell(z,\theta),
\]
where for SFT we can use the standard token-level negative log-likelihood
\[
\ell((x,y),\theta) := -\log p_\theta(y\,|\,x).
\]

We consider a small upweighting of the synthetic subset $S$:
\[
\mathcal{L}_\varepsilon(\theta) := \mathcal{L}_{0}(\theta) + \varepsilon \,\mathcal{L}_{S}(\theta),
\]
where $\mathcal{L}_{0}$ is a \emph{reference} training objective defining the reference point $\theta^\star$.
Pragmatically, $\theta^\star$ can be (i) a checkpoint before adding the new synthetic data,
or (ii) a solution trained on a broader mixture so that group-wise gradients need not vanish exactly.
This avoids degeneracy where $\nabla \mathcal{L}_{\mathrm{real}}(\theta^\star)=0$ would remove real--synthetic interaction signals.

Let
\[
\theta^\star := \arg\min_\theta \mathcal{L}_{0}(\theta),
\qquad
\theta_\varepsilon := \arg\min_\theta \mathcal{L}_{\varepsilon}(\theta).
\]

\subsection{Evaluation function $f$ for LLM tuning}
We choose a differentiable evaluation objective $f(\theta)$ to \emph{minimize}, defined on held-out evaluation sets.
A practical choice is a weighted sum of differentiable surrogate losses:
\[
f(\theta) = \lambda_{\mathrm{ID}} f_{\mathrm{ID}}(\theta) + \lambda_{\mathrm{OOD}} f_{\mathrm{OOD}}(\theta)
+ \lambda_{\mathrm{safe}} f_{\mathrm{safe}}(\theta) + \lambda_{\mathrm{pref}} f_{\mathrm{pref}}(\theta),
\]
where:
\begin{align*}
f_{\mathrm{ID}}(\theta)
&:= \frac{1}{|\mathcal{D}_{\mathrm{eval}}^{\mathrm{ID}}|}\sum_{(x,y)\in \mathcal{D}_{\mathrm{eval}}^{\mathrm{ID}}} \big(-\log p_\theta(y|x)\big),\\
f_{\mathrm{OOD}}(\theta)
&:= \frac{1}{|\mathcal{D}_{\mathrm{eval}}^{\mathrm{OOD}}|}\sum_{(x,y)\in \mathcal{D}_{\mathrm{eval}}^{\mathrm{OOD}}} \big(-\log p_\theta(y|x)\big),\\
f_{\mathrm{safe}}(\theta)
&:= \frac{1}{|\mathcal{D}_{\mathrm{eval}}^{\mathrm{safe}}|}\sum_{(x,y_{\mathrm{safe}})\in \mathcal{D}_{\mathrm{eval}}^{\mathrm{safe}}} \big(-\log p_\theta(y_{\mathrm{safe}}|x)\big).
\end{align*}
If preference data is available, a differentiable held-out preference objective (e.g., DPO-style) can be used:
\[
f_{\mathrm{pref}}(\theta)
:= \frac{1}{|\mathcal{D}_{\mathrm{eval}}^{\mathrm{pref}}|}\sum_{(x,y^+,y^-)\in \mathcal{D}_{\mathrm{eval}}^{\mathrm{pref}}}
\Big(-\log \sigma\big(\beta(\log p_\theta(y^+|x)-\log p_\theta(y^-|x))\big)\Big),
\]
with inverse-temperature $\beta>0$ and sigmoid $\sigma(\cdot)$.
The weights $\lambda_\cdot$ set the trade-off between in-domain quality, generalization, safety, and preferences.

\subsection{Second-order influence for adding a synthetic subset}
Define at $\theta^\star$:
\[
g_f := \nabla_\theta f(\theta^\star), \quad H_f := \nabla_\theta^2 f(\theta^\star), \quad
H := \nabla_\theta^2 \mathcal{L}_0(\theta^\star),
\]
and for the selected synthetic subset $S$,
\[
g_S := \nabla_\theta \mathcal{L}_{S}(\theta^\star), \quad H_S := \nabla_\theta^2 \mathcal{L}_{S}(\theta^\star).
\]

The second-order influence expansion gives
\begin{equation}
f(\theta_\varepsilon)
\approx
f(\theta^\star)
-\varepsilon \, g_f^\top H^{-1} g_S
+\frac{\varepsilon^2}{2}
\Big(
(H^{-1}g_S)^\top H_f (H^{-1}g_S)
+2(H^{-1}g_f)^\top H_S (H^{-1}g_S)
\Big).
\label{eq:second_order_if}
\end{equation}
In what follows, we focus on the synergy term
\[
\Psi(S) := (H^{-1}g_S)^\top H_f (H^{-1}g_S),
\]
which captures subset interactions and admits a pairwise expansion.

\subsection{Pairwise interaction view and complementarity}
Let $v(z) := H^{-1}\nabla_\theta \ell(z,\theta^\star)$.
For $S=\{z_1,\dots,z_m\}$ and $g_S=\frac{1}{m}\sum_{i=1}^m \nabla_\theta \ell(z_i,\theta^\star)$,
\[
H^{-1}g_S = \frac{1}{m}\sum_{i=1}^m v(z_i),
\qquad
\Psi(S)
=
\Big(\frac{1}{m}\sum_{i=1}^m v(z_i)\Big)^\top
H_f
\Big(\frac{1}{m}\sum_{j=1}^m v(z_j)\Big)
=
\frac{1}{m^2}\sum_{i,j=1}^m v(z_i)^\top H_f v(z_j).
\]
Thus, the synergy is governed by pairwise similarities under the curvature induced by $H_f$.

\paragraph{Real--synthetic complementarity via a cross term.}
Let $R$ denote a fixed real anchor group (e.g., $\mathcal{D}_{\mathrm{real}}$ or a real minibatch distribution),
with
\[
g_R := \nabla_\theta \mathcal{L}_{R}(\theta^\star), \qquad
\text{and}\qquad
g_{R\cup S} = g_R + g_S.
\]
Then the synergy decomposes as
\begin{equation}
\Psi(R\cup S)
=
(H^{-1}g_R)^\top H_f (H^{-1}g_R)
+
(H^{-1}g_S)^\top H_f (H^{-1}g_S)
+
2(H^{-1}g_R)^\top H_f (H^{-1}g_S).
\label{eq:cross_decomp}
\end{equation}
The cross term
\[
\Psi_{\times}(R,S) := (H^{-1}g_R)^\top H_f (H^{-1}g_S)
\]
quantifies how synthetic updates interact with real updates under the evaluation curvature.

\paragraph{Sign convention (minimization).}
Because $f$ is defined as an evaluation \emph{loss} to be minimized, the net predicted gain from adding $S$
includes a term $-\varepsilon^2\,\Psi_{\times}(R,S)$ (see below).
Hence, for loss-minimization, a \emph{negative} $\Psi_{\times}(R,S)$ is constructive (complementary),
while a positive $\Psi_{\times}(R,S)$ is potentially interfering.

\subsection{Objective for selecting complementary synthetic subsets}
We compare training with real anchor $R$ alone versus $R\cup S$.
Ignoring the $H_S$-dependent correction in \eqref{eq:second_order_if} for clarity (or treating it as a higher-cost optional term),
the predicted change in evaluation loss from adding $S$ on top of $R$ is
\begin{align}
\widehat{\Delta f}(S;R)
&:= \widehat{f}(R\cup S) - \widehat{f}(R) \nonumber\\
&\approx
-\varepsilon\, g_f^\top H^{-1} g_S
+\frac{\varepsilon^2}{2}\Big(
2(H^{-1}g_R)^\top H_f(H^{-1}g_S)
+
(H^{-1}g_S)^\top H_f(H^{-1}g_S)
\Big).
\label{eq:delta_f_RS}
\end{align}
Define the predicted \emph{gain} (decrease in $f$)
\[
\widehat{\mathcal{G}}(S;R) := -\widehat{\Delta f}(S;R).
\]
Then
\begin{equation}
\widehat{\mathcal{G}}(S;R)
\approx
\underbrace{\varepsilon\, g_f^\top H^{-1} g_S}_{\text{first-order utility}}
\;-\;
\underbrace{\varepsilon^2 (H^{-1}g_R)^\top H_f(H^{-1}g_S)}_{\text{real--synthetic complementarity}}
\;-\;
\underbrace{\frac{\varepsilon^2}{2}(H^{-1}g_S)^\top H_f(H^{-1}g_S)}_{\text{synthetic redundancy / diminishing returns}}.
\label{eq:gain_objective}
\end{equation}
We select $S$ by maximizing $\widehat{\mathcal{G}}(S;R)$ subject to a budget:
\[
\max_{S\subseteq \mathcal{P}_{\mathrm{syn}}} \;\widehat{\mathcal{G}}(S;R)
\quad
\text{s.t.}\quad |S|\le B \;\; \text{or}\;\; \mathrm{Tok}(S)\le T.
\]

\paragraph{Per-example decomposition.}
Let $S=\{z_{i_1},\dots,z_{i_m}\}$ and define $v_i := v(z_i)=H^{-1}\nabla \ell(z_i,\theta^\star)$.
Let
\[
u_i := g_f^\top v_i, \qquad c_i := (H^{-1}g_R)^\top H_f v_i, \qquad q_{ij} := v_i^\top H_f v_j.
\]
Then
\[
\widehat{\mathcal{G}}(S;R)
\approx
\frac{\varepsilon}{m}\sum_{i\in S} u_i
-\frac{\varepsilon^2}{m}\sum_{i\in S} c_i
-\frac{\varepsilon^2}{2m^2}\sum_{i,j\in S} q_{ij}.
\]
The last term penalizes choosing many synthetic points that are similar under $q_{ij}$,
which operationalizes a curvature-aware notion of redundancy in the Hessian-weighted gradient space.

\paragraph{Greedy selection (practical).}
A practical approach is greedy maximization using marginal gains:
for current set $S$, add $z_i$ that maximizes $\Delta(i\mid S) := \widehat{\mathcal{G}}(S\cup\{z_i\};R)-\widehat{\mathcal{G}}(S;R)$.
This yields an interaction-aware data selection rule that prefers
high first-order utility, negative cross interaction with $R$ (constructive complementarity for loss-minimization),
and low redundancy with the already selected synthetic points.

\subsection{Projected computation for LLM-scale models}
Directly forming $H^{-1}$ and $H_f$ is infeasible for LLMs.
Introduce a projection $P\in\mathbb{R}^{k\times n}$ with $k\ll n$ and define
\[
\tilde{H} := P H P^\top,\quad \tilde{H}_f := P H_f P^\top,\quad \tilde{g}(z) := P\nabla_\theta \ell(z,\theta^\star),\quad \tilde{g}_f := P g_f,\quad \tilde{g}_R := P g_R.
\]
Approximate
\[
H^{-1}g \approx P^\top \tilde{H}^{-1}(P g),\qquad
v(z)=H^{-1}\nabla \ell(z,\theta^\star)\approx P^\top \tilde{H}^{-1}\tilde{g}(z).
\]
Then the key quantities become low-dimensional:
\[
u_i \approx \tilde{g}_f^\top \tilde{H}^{-1}\tilde{g}(z_i),
\qquad
c_i \approx (\tilde{H}^{-1}\tilde{g}_R)^\top \tilde{H}_f(\tilde{H}^{-1}\tilde{g}(z_i)),
\qquad
q_{ij} \approx (\tilde{H}^{-1}\tilde{g}(z_i))^\top \tilde{H}_f(\tilde{H}^{-1}\tilde{g}(z_j)).
\]
All computations scale with $k$ rather than $n$.

\subsection{What to expect (testable hypotheses)}
Assuming the approximation is sufficiently accurate for ranking, we expect:
\begin{enumerate}
  \item \textbf{Token efficiency:} at a fixed synthetic token budget, the selected set $S$ yields lower evaluation loss
  (or better downstream metrics) than baselines such as random selection and first-order influence-only selection.
  \item \textbf{Complementarity over redundancy:} selected synthetic examples exhibit more negative cross interaction
  $\Psi_{\times}(R,S)$ (constructive for loss-minimization) compared to naive diversity heuristics,
  and lower within-synthetic redundancy measured by the quadratic term.
  \item \textbf{Improved generalization:} when $f$ includes an OOD component, selection preferentially chooses synthetic data
  that improves OOD while minimally degrading in-domain and safety components.
  \item \textbf{Predictive validity:} the predicted gain $\widehat{\mathcal{G}}(S;R)$ correlates with empirical improvements
  measured after fine-tuning (rank correlation across different subsets).
\end{enumerate}

\subsection{Experimental plan}
\paragraph{Step 0: model and training regime.}
Choose a base LLM and a tuning regime (full fine-tuning or parameter-efficient tuning such as LoRA).
Fix the real dataset $\mathcal{D}_{\mathrm{real}}$ and prepare a large synthetic pool $\mathcal{P}_{\mathrm{syn}}$.

\paragraph{Step 1: define evaluation objective $f$.}
Construct evaluation sets:
$\mathcal{D}_{\mathrm{eval}}^{\mathrm{ID}}$ (held-out real),
$\mathcal{D}_{\mathrm{eval}}^{\mathrm{OOD}}$ (gap-focused tasks),
$\mathcal{D}_{\mathrm{eval}}^{\mathrm{safe}}$ (safety),
and optionally $\mathcal{D}_{\mathrm{eval}}^{\mathrm{pref}}$ (preferences).
Define $f$ as the weighted sum in the previous section.

\paragraph{Step 2: choose reference point $\theta^\star$.}
Use a checkpoint that reflects the current model state before adding the candidate synthetic subset.
Compute (or approximate) $H$, $g_f$, $H_f$, and the group gradient $g_R$ at $\theta^\star$.

\paragraph{Step 3: compute projected quantities.}
Choose projection dimension $k$ and projection matrix $P$.
Compute $\tilde{H}$ and $\tilde{H}_f$ using stochastic approximations (e.g., Gauss--Newton / Fisher-style estimates).
For each candidate synthetic example $z_i\in\mathcal{P}_{\mathrm{syn}}$, compute $\tilde{g}(z_i)$.

\paragraph{Step 4: selection methods (compare baselines).}
For each budget $B$ (or token budget $T$), construct subsets:
\begin{itemize}
  \item \textbf{Random:} uniform sampling from $\mathcal{P}_{\mathrm{syn}}$.
  \item \textbf{First-order only:} select top-$B$ by $u_i$ (utility) or by $g_f^\top H^{-1} g_{z_i}$.
  \item \textbf{Heuristic diversity:} cluster-based uniform sampling (embedding-space or gradient-space).
  \item \textbf{Second-order complement (proposed):} greedy maximize $\widehat{\mathcal{G}}(S;R)$ in \eqref{eq:gain_objective}
  using projected computations.
  \item \textbf{Ablations:} (i) remove cross term; (ii) remove redundancy term; (iii) vary $k$.
\end{itemize}

\paragraph{Step 5: controlled fine-tuning.}
For each selected subset $S$, fine-tune from the same initialization with the same compute budget.
Keep the real data proportion fixed and add synthetic subset $S$ according to a controlled schedule.

\paragraph{Step 6: evaluation and analysis.}
Report:
\begin{itemize}
  \item \textbf{Primary:} evaluation losses corresponding to $f_{\mathrm{ID}}, f_{\mathrm{OOD}}, f_{\mathrm{safe}}, f_{\mathrm{pref}}$.
  \item \textbf{Secondary:} non-differentiable downstream metrics (benchmark scores, judge win-rates, red-teaming success rates).
  \item \textbf{Attribution validity:} rank correlation between $\widehat{\mathcal{G}}(S;R)$ and empirical gains.
  \item \textbf{Interaction diagnostics:} distributions of $\Psi_{\times}(R,S)$ and within-synthetic redundancy for each method.
\end{itemize}

\paragraph{Step 7: robustness checks.}
Repeat across:
(i) different synthetic generators,
(ii) different real dataset sizes (to test complementarity under scarcity),
(iii) different evaluation weightings $\lambda_\cdot$,
(iv) multiple random seeds for selection and training.


% ============================================================
% Influence estimation for non-differentiable evaluation metrics
% ============================================================

\section{Influence Estimation with Non-differentiable Evaluation Metrics}
\label{sec:nondiff_f}

\subsection{Why non-differentiable $f$ arises in LLM evaluation}
Many LLM evaluations are \emph{discrete} or \emph{procedural}:
exact match, pass@k, unit-test success, toxicity thresholding,
and win-rate under an external judge (human or LLM-as-judge).
Such metrics are non-differentiable with respect to model parameters $\theta$,
especially under deterministic decoding.
However, the second-order influence expansion used in this work requires
$g_f=\nabla_\theta f(\theta^\star)$ and $H_f=\nabla_\theta^2 f(\theta^\star)$
in Eq.~(1).
This section provides two ways to handle non-differentiable $f$:

\begin{itemize}
\item \textbf{Stochastic relaxation (differentiable in expectation):}
define $f$ as the expected value of a non-differentiable metric under stochastic generation,
then compute unbiased estimators of $g_f$ and (if needed) $H_f$ via score-function identities.
\item \textbf{Black-box directional estimation:}
avoid explicit $g_f$ and $H_f$ by estimating the directional derivatives of $f$
along influence directions using a small number of function evaluations.
\end{itemize}

\subsection{Stochastic relaxation: expected metric objective}
Let $\mathcal{D}_{\mathrm{eval}}=\{x_i\}_{i=1}^M$ be evaluation prompts.
Let the model induce a conditional distribution over outputs
$y \sim p_\theta(\cdot \mid x)$ (e.g., sampling at temperature $\tau>0$).
Let $r(x,y)\in \mathbb{R}$ be a (possibly non-differentiable) reward computed from a metric or a judge,
for example:
\begin{itemize}
\item $r(x,y)=\mathbf{1}\{\text{unit tests pass}\}$ for code,
\item $r(x,y)=\mathbf{1}\{\text{exact match}\}$ for QA,
\item $r(x,y)=\mathbf{1}\{\text{judge prefers } y \text{ over baseline}\}$ for win-rate,
\item $r(x,y)\in\{0,1,2\}$ for graded judge outcomes.
\end{itemize}
Define the evaluation \emph{loss} to minimize as the negative expected reward:
\begin{equation}
f(\theta)
:=
-\frac{1}{M}\sum_{i=1}^M
\mathbb{E}_{y \sim p_\theta(\cdot\mid x_i)}\big[r(x_i,y)\big].
\label{eq:expected_metric_f}
\end{equation}
Even if $r(\cdot,\cdot)$ is non-differentiable, $f(\theta)$ is differentiable under mild regularity conditions,
because $\theta$ only enters through $p_\theta(y\mid x)$.

\paragraph{First derivative (score-function / REINFORCE).}
Let $s_\theta(x,y):=\nabla_\theta \log p_\theta(y\mid x)$.
Then
\begin{equation}
g_f(\theta)
=
\nabla_\theta f(\theta)
=
-\frac{1}{M}\sum_{i=1}^M
\mathbb{E}_{y \sim p_\theta(\cdot\mid x_i)}
\Big[r(x_i,y)\, s_\theta(x_i,y)\Big].
\label{eq:gf_reinforce}
\end{equation}
A standard variance-reduced form introduces a baseline $b(x)$:
\[
g_f(\theta)
=
-\frac{1}{M}\sum_{i=1}^M
\mathbb{E}_{y \sim p_\theta(\cdot\mid x_i)}
\Big[(r(x_i,y)-b(x_i))\, s_\theta(x_i,y)\Big],
\]
since $\mathbb{E}[b(x)\,s_\theta(x,y)]=b(x)\nabla_\theta \sum_y p_\theta(y\mid x)=0$.

\paragraph{Second derivative (Hessian identity).}
Assuming $r(x,y)$ does not depend on $\theta$ (judge fixed),
\begin{align}
H_f(\theta)
&=
\nabla_\theta^2 f(\theta)\nonumber\\
&=
-\frac{1}{M}\sum_{i=1}^M
\mathbb{E}_{y \sim p_\theta(\cdot\mid x_i)}
\Big[
r(x_i,y)\,\nabla_\theta^2 \log p_\theta(y\mid x_i)
+
r(x_i,y)\, s_\theta(x_i,y)s_\theta(x_i,y)^\top
\Big].
\label{eq:Hf_score_function}
\end{align}
In practice, one rarely forms $H_f$ explicitly; one estimates Hessian-vector products
or only the quadratic forms $u^\top H_f u$ needed by second-order influence.

\subsection{Plug-in influence with expected metric $f$}
Let $\theta^\star$ be the reference parameters and let $H=\nabla_\theta^2 \mathcal{L}_0(\theta^\star)$.
For a training subset $S$ with $g_S=\nabla_\theta \mathcal{L}_S(\theta^\star)$ and $H_S=\nabla_\theta^2 \mathcal{L}_S(\theta^\star)$,
the second-order influence approximation (Eq.~(1)) applies with $f$ defined by \eqref{eq:expected_metric_f}:
\begin{equation}
f(\theta_\varepsilon)
\approx
f(\theta^\star)
-\varepsilon\, g_f^\top H^{-1}g_S
+\frac{\varepsilon^2}{2}
\Big(
(H^{-1}g_S)^\top H_f (H^{-1}g_S)
+2(H^{-1}g_f)^\top H_S(H^{-1}g_S)
\Big),
\label{eq:second_order_if_nondiff}
\end{equation}
where $g_f$ and $H_f$ can be estimated via \eqref{eq:gf_reinforce}--\eqref{eq:Hf_score_function} at $\theta=\theta^\star$.

\paragraph{Quadratic form of $H_f$ without constructing $H_f$.}
Let $u \in \mathbb{R}^n$. From \eqref{eq:Hf_score_function},
\begin{equation}
u^\top H_f(\theta) u
=
-\frac{1}{M}\sum_{i=1}^M
\mathbb{E}_{y \sim p_\theta(\cdot\mid x_i)}
\Big[
r(x_i,y)\, u^\top \nabla_\theta^2 \log p_\theta(y\mid x_i)\, u
+
r(x_i,y)\, (u^\top s_\theta(x_i,y))^2
\Big].
\label{eq:ufHu_score}
\end{equation}
Thus the synergy term in \eqref{eq:second_order_if_nondiff} can be computed using only:
(i) scalar products $u^\top s_\theta$, and
(ii) Hessian-vector products of $\nabla_\theta^2 \log p_\theta(y\mid x)$ with $u$.

\subsection{Black-box influence estimation via directional finite differences}
If the evaluation metric is available only as a black-box function
(e.g., proprietary judge, human-in-the-loop, complex simulator),
we can avoid explicit $g_f$ and $H_f$ by estimating the required directional derivatives.

\paragraph{Influence direction.}
Define the (first-order) influence direction for subset $S$:
\begin{equation}
d_S := -H^{-1}g_S.
\label{eq:influence_direction}
\end{equation}
This is the parameter change induced by upweighting $S$ by $\varepsilon$ to first order.
(For LLMs, $d_S$ can be approximated in a low-dimensional subspace via gradient projection.)

\paragraph{First-order influence as a directional derivative.}
Assuming $f$ is directionally differentiable at $\theta^\star$ along $d_S$,
\begin{equation}
-\varepsilon\, g_f^\top H^{-1}g_S
=
\varepsilon\, g_f^\top d_S
=
\varepsilon\,\frac{\mathrm{d}}{\mathrm{d}t} f(\theta^\star + t d_S)\Big|_{t=0}.
\label{eq:first_order_directional}
\end{equation}
We can estimate the derivative using a one-sided finite difference:
\begin{equation}
\widehat{\mathcal{I}}^{(1)}(S)
:=
\varepsilon \cdot \frac{f(\theta^\star + h d_S)-f(\theta^\star)}{h},
\label{eq:first_order_fd}
\end{equation}
with step size $h>0$.

\paragraph{Second-order synergy as a second directional derivative.}
The synergy term equals the second directional derivative of $f$ along $d_S$:
\begin{equation}
(H^{-1}g_S)^\top H_f (H^{-1}g_S)
=
d_S^\top H_f d_S
=
\frac{\mathrm{d}^2}{\mathrm{d}t^2} f(\theta^\star + t d_S)\Big|_{t=0}.
\label{eq:second_order_directional}
\end{equation}
A centered finite difference yields:
\begin{equation}
\widehat{\Psi}(S)
:=
\frac{f(\theta^\star + h d_S) -2f(\theta^\star)+ f(\theta^\star - h d_S)}{h^2}.
\label{eq:second_order_fd}
\end{equation}
Then the second-order contribution in Eq.~(1) can be approximated as
$\frac{\varepsilon^2}{2}\widehat{\Psi}(S)$ (ignoring the $H_S$ correction term),
or used directly as an interaction-aware selection score.

\paragraph{Cross interaction between real anchor $R$ and synthetic subset $S$.}
For complementarity analysis, the cross term is the mixed second derivative:
\begin{equation}
(H^{-1}g_R)^\top H_f (H^{-1}g_S)
=
d_R^\top H_f d_S
=
\frac{\partial^2}{\partial \alpha\,\partial \beta}
f(\theta^\star + \alpha d_R + \beta d_S)\Big|_{\alpha=\beta=0}.
\label{eq:mixed_second_derivative}
\end{equation}
It can be estimated via a four-point stencil:
\begin{equation}
\widehat{\Psi}_\times(R,S)
:=
\frac{
f(\theta^\star + h d_R + h d_S)
- f(\theta^\star + h d_R)
- f(\theta^\star + h d_S)
+ f(\theta^\star)
}{h^2}.
\label{eq:cross_fd}
\end{equation}
This estimator uses only evaluations of $f$, not its gradients.

\subsection{Practical notes (variance and feasibility)}
\paragraph{Common random numbers.}
When $f(\theta)$ is estimated by sampling generations and judging them,
use the same random seeds (or the same sampled noise) across
$\theta^\star$, $\theta^\star \pm h d_S$, and $\theta^\star + h(d_R+d_S)$.
This coupling greatly reduces the variance of finite differences.

\paragraph{Where to perturb parameters.}
For large LLMs, apply perturbations only in a tractable subspace:
adapter weights (LoRA) or a projected subspace $P^\top \mathbb{R}^k$.
Then $d_S$ in \eqref{eq:influence_direction} is replaced by a projected direction,
and $\theta^\star \pm h d_S$ can be realized by a low-rank parameter delta.

\paragraph{Step-size selection.}
Pick $h$ such that:
(i) the metric estimator is stable (not dominated by Monte Carlo noise),
and (ii) the Taylor approximation remains locally valid.
A practical procedure is to sweep $h$ over a logarithmic grid
and choose a region where the finite-difference estimate is stable.

\paragraph{Summary.}
Non-differentiable evaluations can be handled either by:
(i) defining an expected-metric $f$ and using score-function estimators for $g_f,H_f$,
or (ii) treating $f$ as black-box and estimating the first-order and second-order influence terms
from a small number of metric evaluations along influence directions.
\end{document}
